\SourceLeaf{392}%
\Needspace{9\baselineskip}
\begin{center}\bfseries Remarks.\end{center}
\noindent\hspace*{1.2em} On 11 December 1822, in answer to a prize question posed by the Royal Society of Sciences in Copenhagen, \textsc{Gauss} had sent his paper \textit{Allgemeine Auflösung der Aufgabe: die Theile einer gegebenen Fläche auf einer andern gegebenen Fläche so abzubilden, dass die Abbildung dem Abgebildeten in den kleinsten Theilen ähnlich wird} (vol. IV, pp. 189 to 216) to \textsc{Schumacher}. The preceding notes from 13 to 15 December 1822, written on a single leaf, show that he had quite deliberately sent it out with \textsc{Newton}'s proud motto »Ab his via sternitur ad maiora«. For the representation of the square of the line element of an arbitrary surface in the form
\[
mm(dt^2+du^2)
\]
had opened the way to solving the problem of expressing the measure of curvature, whose invariance under bending he had already recognized early on (cf. p. 372), in terms of the coefficients of the line element, here $mm$, and their derivatives. The foundation for the \textit{Disquisitiones generales circa superficies curvas} was thereby laid.

A note in a notebook concerns the same subject and presumably dates from about the same time as the entry already printed on p. 373.

{\itshape »The following very elegant general theorem probably holds:

If $n=e^\nu$, then
\[
\begin{aligned}
\frac{\partial\partial\nu}{\partial x^2}+\frac{\partial\partial\nu}{\partial y^2}
&=\frac{n\left(\frac{\partial\partial n}{\partial x^2}+\frac{\partial\partial n}{\partial y^2}\right)-\left\{\left(\frac{\partial n}{\partial x}\right)^2+\left(\frac{\partial n}{\partial y}\right)^2\right\}}{nn}\\
&=\frac1{rr'nn}=\frac{d\Sigma}{dS.nn}=\frac{d\Sigma}{ds},
\end{aligned}
\]
if the two extreme radii of curvature, perpendicular to one another, are what $r$ and $r'$

\noindent [denote and] if
\begin{center}
\begin{tabular}{ccccc}
$d\Sigma$&area&on&the&celestial sphere\\
$dS$&,,&,,&,,&curved surface\\
$ds$&,,&\multicolumn{3}{l}{in the plane}
\end{tabular}
\end{center}
[and $ds=nn\,dS$ holds.]«\par}
\begin{flushright}\textsc{Stäckel.}\end{flushright}
